\documentclass[11pt,a4paper]{article}
\usepackage[utf8]{inputenc}
\usepackage[T1]{fontenc}
\usepackage[margin=1in]{geometry}
\usepackage{amsmath,amssymb}
\usepackage{listings}
\usepackage{xcolor}
\usepackage{hyperref}

\lstset{
  language=C++,
  basicstyle=\ttfamily\small,
  keywordstyle=\color{blue}\bfseries,
  commentstyle=\color{green!50!black},
  stringstyle=\color{red!70!black},
  numbers=left,
  numberstyle=\tiny\color{gray},
  breaklines=true,
  frame=single,
  tabsize=4,
  showstringspaces=false
}

\title{IOI 2012 -- Parrots}
\author{Solution}
\date{}

\begin{document}
\maketitle

\section{Problem Summary}
Encode a message of $N \le 64$ bytes into a multiset of integers in $[0,255]$. The decoder receives the integers in arbitrary order and must reconstruct the original message exactly.

\section{Solution}

\subsection{Messages as Ranks}
Interpret the byte array as a base-$256$ integer
\[
V = \sum_{i=0}^{N-1} M_i \cdot 256^i .
\]
There are exactly $256^N$ possible messages.

\subsection{Encoded Multisets as Weak Compositions}
Fix a length bound $T$. Any encoded multiset of length at most $T$ can be written as counts
\[
c_0, c_1, \ldots, c_{255}, c_{256},
\]
where $c_x$ is the number of occurrences of value $x$, and $c_{256}$ is a ``blank'' count so that
\[
\sum_{x=0}^{256} c_x = T.
\]
Thus the number of possible encodings is the number of weak compositions of $T$ into $257$ parts:
\[
\binom{T + 256}{256}.
\]
Choose the smallest $T$ such that $\binom{T + 256}{256} \ge 256^N$. For $N = 64$, this gives $T = 261$, which is the optimal full-score bound.

\subsection{Ranking and Unranking}
We map each message rank $V$ to the $V$-th weak composition in lexicographic order.

Suppose we still have $R$ slots to distribute and we are deciding the count of the current value. If we set it to $t$, then the number of completions is
\[
\binom{(R - t) + K - 2}{K - 2},
\]
where $K$ is the number of categories still available (current value, larger values, and blanks).

This gives:
\begin{itemize}
  \item an \emph{unranking} procedure for the encoder, which converts $V$ into counts $c_x$;
  \item a \emph{ranking} procedure for the decoder, which reconstructs $V$ from the received counts.
\end{itemize}

\subsection{Why the Order Does Not Matter}
The sent data is only the multiset: value $x$ is transmitted exactly $c_x$ times. Since the decoder only needs the multiplicities, permutation of the transmitted numbers is harmless.

\section{Implementation}

The code uses a tiny fixed-size big integer (10 limbs of 64 bits), which is enough because all values fit comfortably below $2^{640}$.

\begin{lstlisting}
#include <bits/stdc++.h>
using namespace std;

void send(int value);
void output(int value);

namespace {
constexpr int K = 257;
constexpr int MAX_N = 64;
constexpr int MAX_T = 261;
constexpr int MAX_C = 256 + MAX_T;
constexpr int LIMBS = 10;

struct BigInt {
    array<uint64_t, LIMBS> limb{};
    BigInt(uint64_t value = 0) {
        limb.fill(0);
        limb[0] = value;
    }
    bool operator<(const BigInt& other) const {
        for (int i = LIMBS - 1; i >= 0; --i)
            if (limb[i] != other.limb[i])
                return limb[i] < other.limb[i];
        return false;
    }
    BigInt& operator+=(const BigInt& other) {
        unsigned __int128 carry = 0;
        for (int i = 0; i < LIMBS; ++i) {
            unsigned __int128 sum =
                (unsigned __int128)limb[i] + other.limb[i] + carry;
            limb[i] = (uint64_t)sum;
            carry = sum >> 64;
        }
        return *this;
    }
    BigInt& operator-=(const BigInt& other) {
        uint64_t borrow = 0;
        for (int i = 0; i < LIMBS; ++i) {
            uint64_t old = limb[i];
            uint64_t sub = other.limb[i] + borrow;
            limb[i] = old - sub;
            borrow = (old < sub) || (borrow && sub == 0);
        }
        return *this;
    }
    void add_small(uint64_t value) {
        unsigned __int128 carry = value;
        for (int i = 0; i < LIMBS && carry > 0; ++i) {
            unsigned __int128 sum = (unsigned __int128)limb[i] + carry;
            limb[i] = (uint64_t)sum;
            carry = sum >> 64;
        }
    }
    void shift_left_8() {
        uint64_t carry = 0;
        for (int i = 0; i < LIMBS; ++i) {
            uint64_t next = limb[i] >> 56;
            limb[i] = (limb[i] << 8) | carry;
            carry = next;
        }
    }
    void shift_right_8() {
        uint64_t carry = 0;
        for (int i = LIMBS - 1; i >= 0; --i) {
            uint64_t next = limb[i] << 56;
            limb[i] = (limb[i] >> 8) | carry;
            carry = next;
        }
    }
    int low_byte() const { return limb[0] & 255; }
};

vector<vector<BigInt>> binom;
array<int, MAX_N + 1> best_len{};
bool ready = false;

void init_tables() {
    if (ready) return;
    ready = true;

    binom.assign(MAX_C + 1, vector<BigInt>(257, BigInt(0)));
    binom[0][0] = 1;
    for (int n = 1; n <= MAX_C; ++n) {
        binom[n][0] = 1;
        for (int k = 1; k <= min(n, 256); ++k)
            binom[n][k] = binom[n - 1][k - 1];
            binom[n][k] += binom[n - 1][k];
    }

    BigInt states = 1;
    int t = 0;
    for (int n = 1; n <= MAX_N; ++n) {
        states.shift_left_8();
        while (binom[t + 256][256] < states) ++t;
        best_len[n] = t;
    }
}

BigInt message_to_rank(int N, int M[]) {
    BigInt value = 0;
    for (int i = N - 1; i >= 0; --i) {
        value.shift_left_8();
        value.add_small(M[i]);
    }
    return value;
}

vector<int> unrank_counts(int total, BigInt rank) {
    vector<int> cnt(K, 0);
    int rem = total;
    for (int x = 0; x + 1 < K; ++x) {
        int tail = K - x - 2;
        for (int take = 0; take <= rem; ++take) {
            const BigInt& ways = binom[rem - take + tail][tail];
            if (rank < ways) {
                cnt[x] = take;
                rem -= take;
                break;
            }
            rank -= ways;
        }
    }
    cnt.back() = rem;
    return cnt;
}

BigInt rank_counts(const vector<int>& cnt) {
    BigInt rank = 0;
    int rem = accumulate(cnt.begin(), cnt.end(), 0);
    for (int x = 0; x + 1 < K; ++x) {
        int tail = K - x - 2;
        for (int take = 0; take < cnt[x]; ++take)
            rank += binom[rem - take + tail][tail];
        rem -= cnt[x];
    }
    return rank;
}
}  // namespace

void encode(int N, int M[]) {
    init_tables();
    vector<int> cnt = unrank_counts(best_len[N], message_to_rank(N, M));
    for (int x = 0; x < 256; ++x)
        for (int rep = 0; rep < cnt[x]; ++rep)
            send(x);
}

void decode(int N, int L, int X[]) {
    init_tables();
    vector<int> cnt(K, 0);
    for (int i = 0; i < L; ++i) ++cnt[X[i]];
    cnt.back() = best_len[N] - L;

    BigInt rank = rank_counts(cnt);
    for (int i = 0; i < N; ++i) {
        output(rank.low_byte());
        rank.shift_right_8();
    }
}
\end{lstlisting}

\section{Complexity}
\begin{itemize}
  \item Precomputation of binomial coefficients: $O(256 \cdot T)$ states.
  \item Encoding and decoding: $O(256 \cdot T)$ big-integer operations.
  \item Memory: $O(256 \cdot T)$.
\end{itemize}

\end{document}
